\section{Giới thiệu}

\textbf{Thiết kế kiến trúc phần mềm (Software Architecture Design)} là một giai đoạn quan trọng trong vòng đời phát triển phần mềm, đóng vai trò định hình cấu trúc tổng thể của hệ thống thông qua việc xác định các thành phần chính, trách nhiệm, mối quan hệ và cơ chế tương tác giữa chúng. Không chỉ là cơ sở định hướng cho quá trình thiết kế chi tiết và triển khai, kiến trúc phần mềm còn giúp chuyển hóa các yêu cầu nghiệp vụ thành một giải pháp kỹ thuật có tổ chức, nhất quán và phù hợp với mục tiêu của hệ thống. Đặc biệt, một kiến trúc được thiết kế tốt cần đáp ứng đồng thời các yêu cầu chức năng và phi chức năng như hiệu năng, tính bảo mật, độ tin cậy, khả năng mở rộng và khả năng bảo trì. Bên cạnh đó, việc xác định rõ ranh giới thành phần, các nguyên tắc thiết kế và những quyết định kiến trúc quan trọng giúp giảm sự phụ thuộc không cần thiết, hỗ trợ phối hợp giữa các nhóm phát triển, hạn chế rủi ro kỹ thuật và tối ưu chi phí thay đổi trong dài hạn. Vì vậy, thiết kế kiến trúc phần mềm không chỉ là nền tảng cho việc xây dựng một hệ thống đáp ứng yêu cầu hiện tại mà còn quyết định khả năng thích ứng, phát triển và duy trì chất lượng của hệ thống trong tương lai.

Trong khuôn khổ bài tập lớn môn Kiến trúc Phần mềm, nhóm thực hiện đề tài \textbf{Thiết kế kiến trúc cho hệ thống chấm công và kiểm soát ra vào bằng mã QR} được giả định cho Công ty Cổ phần Sản xuất TMC. Hệ thống được xây dựng nhằm hỗ trợ định danh nhân viên thông qua mã QR, kiểm soát quá trình ra vào tại các cổng, ghi nhận dữ liệu chấm công và cung cấp chức năng báo cáo, kiểm toán và tích hợp với hệ thống nhân sự của công ty. Trên cơ sở các yêu cầu nghiệp vụ và những ràng buộc kỹ thuật được xác định trong đề bài, nhóm tập trung phân tích bối cảnh hệ thống, xác định các đặc tính kiến trúc, đánh giá và lựa chọn phong cách kiến trúc phù hợp, xây dựng kiến trúc logic và kiến trúc vật lý, đồng thời xác lập các quyết định và nguyên tắc thiết kế cần thiết. Qua đó, nhóm hướng đến một kiến trúc đáp ứng các yêu cầu về hiệu năng, tính toàn vẹn dữ liệu, bảo mật và khả năng kiểm soát truy cập, làm nền tảng cho quá trình hoàn thiện thiết kế, triển khai và kiểm thử hệ thống ở các giai đoạn tiếp theo.

Bài tập lớn được triển khai theo hai cột mốc chính (Milestone), với phạm vi công việc được xác định theo từng giai đoạn. Cụ thể, \textbf{Milestone 1 (M1)} tập trung vào phân tích yêu cầu và thiết kế kiến trúc phần mềm, bao gồm mô tả bối cảnh hệ thống, xác định các yêu cầu chức năng và phi chức năng, lựa chọn phong cách kiến trúc, xây dựng kiến trúc logic và kiến trúc vật lý, đồng thời xác định các quyết định và nguyên tắc thiết kế. Tiếp theo, \textbf{Milestone 2 (M2)} hướng đến hiện thực hóa phiên bản MVP của hệ thống, đặc tả giao diện người dùng, hợp đồng API và chiến lược kiểm thử nhằm đánh giá mức độ đáp ứng các yêu cầu đã đề ra. Báo cáo này trình bày xuyên suốt cả phần thiết kế kiến trúc (M1) ở các Chương 2--8 và phần hiện thực hệ thống MVP (M2) ở Chương 9.

Báo cáo gồm các chương chính sau:
\begin{itemize}
    \item \textbf{Chương 2 – Nghiên cứu bối cảnh và phân tích nghiệp vụ:} Trình bày bối cảnh doanh nghiệp, mục tiêu và phạm vi hệ thống, các bên liên quan, ranh giới hệ thống, cùng các ràng buộc, giả định và rủi ro.
    \item \textbf{Chương 3 – Thiết kế yêu cầu hệ thống:} Đặc tả các yêu cầu chức năng theo nhóm nghiệp vụ và các yêu cầu phi chức năng, làm cơ sở đánh giá chất lượng và định hướng kiến trúc.
    \item \textbf{Chương 4 – Thiết kế mô hình miền nghiệp vụ:} Mô tả luồng nghiệp vụ chính, các thực thể và quan hệ giữa chúng; trình bày lược đồ lớp ở mức nghiệp vụ.
    \item \textbf{Chương 5 – Đặc tính kiến trúc hệ thống:} Phân tích các đặc tính kiến trúc, lựa chọn ba đặc tính dẫn dắt ưu tiên và đánh giá các đánh đổi.
    \item \textbf{Chương 6 – Cấu trúc kiến trúc hệ thống:} Phân tích architecture quantum, so sánh và lựa chọn architecture style, sau đó trình bày kiến trúc logic và kiến trúc vật lý.
    \item \textbf{Chương 7 – Quyết định kiến trúc hệ thống:} Tổng hợp và phân tích các quyết định kiến trúc về triển khai, xác thực QR, lưu trữ và đồng bộ giao dịch, tích hợp, tổ chức mã nguồn, bảo mật và quyền riêng tư.
    \item \textbf{Chương 8 – Nguyên lý thiết kế hệ thống:} Trình bày các nguyên tắc cấp kiến trúc, cách áp dụng SOLID, cohesion, coupling và tổ chức mã nguồn.
    \item \textbf{Chương 9 – Hiện thực hệ thống:} Trình bày phạm vi MVP, công nghệ và mô hình triển khai container hóa, cấu trúc mã nguồn, thiết kế chi tiết các module cốt lõi, hiện thực chức năng, danh mục màn hình giao diện (sitemap) và đặc tả hợp đồng API.
    \item \textbf{Phụ lục:} Bao gồm các sơ đồ use case và đặc tả bốn use case có ý nghĩa kiến trúc.
\end{itemize}

Nhóm xin chân thành cảm ơn ThS. Lê Đình Thuận đã tận tình giảng dạy, hướng dẫn và hỗ trợ nhóm trong suốt quá trình học tập môn Kiến trúc Phần mềm cũng như thực hiện bài tập lớn. Những kiến thức chuyên môn, định hướng và góp ý quý báu của thầy đã giúp nhóm hiểu rõ hơn về các nguyên lý thiết kế kiến trúc và vận dụng vào quá trình phân tích, xây dựng giải pháp cho hệ thống. Tuy nhiên, do kiến thức và kinh nghiệm thực tiễn còn hạn chế, báo cáo khó tránh khỏi những thiếu sót. Nhóm rất mong nhận được những ý kiến đóng góp từ thầy để hoàn thiện báo cáo và nâng cao kiến thức chuyên môn.